% ================================================================
%  chapters/parts/ch12-b.tex
%  第 12 章  扫描探针显微镜与纳米光谱技术（分片 b）
%  源：S2（下册）第 12 章，书内 p.913—925 / PDF p.399—411
%  范围：12.3（AFM 基本原理：相互作用力、结构与工作模式）、
%        12.4（AFM 数学分析模型）、12.5（AFM 在表面催化相关领域中的应用）
%  不含：12.1—12.2（STM，分片 a）、近场光学 / TERS / Nano-FTIR（分片 c）
%  说明：本分片不写 \chapterstart（分片 a 已写）。
% ================================================================

\subsection{原子力显微镜的基本原理}
\label{sec:afm-principle}

\begin{probout}
原子力显微镜（atomic force microscope，AFM）测量的是\textbf{针尖与样品之间的力}：微悬臂
末端的针尖在样品表面扫描，针尖--样品相互作用力使悬臂发生偏转、或改变其振动状态，检测这
一变化即可反推表面信息。原始输出有三类：$z(x,y)$ 高度像、悬臂振动的振幅/相位（或频移）像，
以及逐点或逐位置的力--距离曲线（力学谱）。它\textbf{不要求样品导电}，因此把原子级表面
成像从金属与半导体扩展到氧化物、分子筛与负载型催化剂；代价是同一幅像里混合了几何形貌、
力学性质与吸附层信息，\textbf{高度不等于几何高度}。
\end{probout}

\subsubsection{与 STM 的根本差别：测力，而不是测电流}

STM 的成像信号是针尖与表面之间的隧穿电流，前提是二者能构成隧穿结：样品必须有一定导电性，
或被减薄到足以导电。这一条把大多数催化剂载体挡在门外——\ce{CeO2}、\ce{TiO2}、\ce{Al2O3}、
分子筛等宽禁带氧化物与绝缘体在常规偏压下给不出稳定的隧穿电流
\src{S2 §12.3，书内 p.913 / PDF p.399}。

AFM 换掉的正是这个前提：它探测针尖尖端原子与样品表面原子之间的\textbf{相互作用力}。
力不需要载流子输运，因此\textbf{导体、半导体与绝缘体一视同仁}——只要针尖能与表面建立
可检测的力，就能成像。这是 AFM 成为氧化物表面结构表征主力工具的原因，也是它在负载型
金属催化剂研究中的立足点：载体是绝缘的，金属纳米粒子长在载体上，只有不依赖导电性的
方法才能在同一次扫描中同时看到二者。反过来，AFM 也因此丢掉了 STM 的一个优势——
STM 电流对局域电子态密度敏感，AFM 的力信号不直接给电子结构，二者互补而非替代
（见 \ref{sec:afm-crosscheck} 段）。

AFM 的技术沿革可以看成一条"把力测得更准"的路线\src{S2 §12.3，书内 p.913 / PDF p.399}：
Binnig 等在早期 STM 实验中注意到针尖与表面之间存在微弱作用力，据此于 1986 年发明第一台
AFM，突破了 STM 对样品导电性的要求；1987 年 Martin 等引入非接触模式，并测量了针尖--样品
不同间距下的力曲线；1991 年基于频率调制的 AFM 问世并成功用于真空中成像；1993 年 Zhong 等
发明轻敲模式，大幅提高成像稳定性，并很快被用于液体中成像；1996 年 Giessibl 把钟表工业
广泛使用的石英音叉（quartz tuning fork）用作 AFM 微悬臂，发明 qPlus 力学传感器，实现
Si(111)-($7\times7$) 表面的原子分辨；2007 年 Custance 等用原子力谱实现合金表面单原子的
元素区分；2009 年 Gross 等用非接触 AFM 实现实空间碳碳化学键的直接观测，把分辨率推进到
亚原子级。目前 AFM 已广泛用于真空、大气与液体中的成像，以及原子操纵和分子识别，其功能
还可扩展到近场光学成像（分片 c）。

\subsubsection{针尖--样品间的相互作用力}

\begin{principle}
AFM 通过探测针尖尖端与样品表面之间的力学信号成像。针尖探测到的总有效相互作用力可取正值
（排斥）或负值（吸引），二者分别定义了排斥与吸引两个成像工作区。总力由几种长程力与短程力
贡献而成：长程力中除静电力外，最常见的是范德华力（van der Waals force），它是吸引性的，
广泛存在于中性原子与分子之间；短程力只在针尖--样品距离很近时才明显，来自针尖尖端原子与
样品表面原子之间的化学成键作用及 Pauli 排斥作用。对不带电的针尖与样品：间距大于
\qty{1}{nm} 时长程范德华作用最大，一般表现为吸引；针尖继续靠近到接近化学成键距离时，
短程吸引性化学力（源于针尖与表面原子最外层电子波函数的交叠）使吸引力达到最大；越过成键
距离后间距继续减小，电子间 Pauli 排斥急剧增加，排斥力占据主导
\src{S2 §12.3.1，书内 p.913--914 / PDF p.399--400}。
\end{principle}

对不带电的针尖与样品，这一距离依赖可以用势能曲线（Lennard-Jones 势）概括：

\begin{equation}\label{eq:lj}
U(d)=4\varepsilon\left[\left(\frac{\sigma}{d}\right)^{12}-\left(\frac{\sigma}{d}\right)^{6}\right]
\end{equation}

式中 $d$ 为针尖--样品间距，$\varepsilon$ 为势阱深度，$\sigma$ 为零势能间距；$d^{-12}$ 项
代表短程 Pauli 排斥，$d^{-6}$ 项代表长程范德华吸引。力为势能的负梯度，
$F_{\mathrm{ts}}=-\dd U/\dd d$。

\tocheck{未对照原页面图像：源书正文只说"势能曲线（Lennard-Jones 势）可以反映针尖--表面
相互作用随间距的变化特点"，未给出解析式；式~\eqref{eq:lj} 按标准 Lennard-Jones 形式补写，
源书对 $\varepsilon$、$\sigma$ 的符号与取值约定未见正文定义。}

源书图 12-21（a）把隧穿电流与针尖--样品作用力随距离的变化画在同一张图上，可直接读出三件事：
隧穿电流随距离减小而单调增加；针尖--样品作用力先由吸引主导并随距离减小而增大到最大吸引，
接近成键距离时排斥力不断增加，以致在总力为零处到达成键距离，成键距离以下则由急剧增加的
排斥作用主导；因此\textbf{若 AFM 能稳定探测短程力，针尖就能在比 STM 更近的距离下操作，
从而获得更好的分辨}\src{S2 §12.3.1，书内 p.914 / PDF p.400}。这句判断值得单独记住：AFM 的
分辨率上限并不来自"力"本身，而来自能否把短程力从长程背景中分离出来——后面三种工作模式的
差别，本质上就是分离方式与工作距离的差别。

\subsubsection{微悬臂探针与偏转检测}

AFM 的主要组成部分与 STM 十分类似：基于压电陶瓷的针尖--样品运动系统、反馈控制电路单元，
以及探针检测系统。AFM 的核心在力学传感器，目前广泛采用的是微悬臂探针
（cantilever--tip）系统：微悬臂前端有突出的针尖，在不同的针尖--样品力学作用下微悬臂发生
相应偏转（deflection），根据偏转程度即可反映表面信息
\src{S2 §12.3.2，书内 p.914 / PDF p.400}。

\begin{measure}
悬臂偏转最常用的检测技术有两类：一是基于光束发射或干涉的\textbf{偏转检测}（光束偏转法），
二是基于压电效应的\textbf{自检测}。光束偏转法把激光聚焦在悬臂背面，反射光落到四象限光电
探测器上，悬臂偏转使光斑在探测器上位移，位移量转成电压信号（图~\ref{fig:afm-setup}）；
它是"光杠杆"，把纳米级的悬臂偏转放大成探测器上毫米级的光斑移动，这是 AFM 能分辨皮米量级
位移的原因。写入正文时真正需要记录的条件只有四组：\textbf{悬臂弹性常数 $k$}（决定同样偏转
对应多大的力，也决定会不会压坏样品）、\textbf{共振频率 $f_0$ 与品质因数 $Q$}（决定动态
模式的振幅稳定性与响应速度）、\textbf{针尖半径 $R$}（决定横向分辨率与卷积加宽，见
\ref{sec:afm-conv} 段）、以及\textbf{环境}（真空/大气/液体，直接决定表面有无吸附水膜）
\src{S2 §12.3.2，书内 p.914--915 / PDF p.400--401}。任何一条没有随图给出，力曲线的定量
结论就不能跨实验比较。
\end{measure}

\begin{figure}[htbp]
\centering
\begin{tikzpicture}[x=1mm,y=1mm,>=Stealth]
% ---- 坐标区间静态核算（x=1mm，y=1mm；版心 165 mm）----
% 图元：样品矩形 (30,10)-(84,17)；xy 扫描管 (44,2)-(70,9)；悬臂支架 (12,36)-(26,44)；
%       悬臂线 y=40（x 24→64）；针尖三角 (61,40)-(67,40)-(64,19)；
%       四象限探测器 (68,58)-(92,72)；反馈盒 (96,32)-(130,46)。
% 节点区间（anchor 已注明，宽度按 text width 计）：
%  N1 样品标签   中心(57,13.5) w48 → x[33,81]  y[11.0,16.0]  在样品矩形内
%  N2 扫描管标签 中心(57,5.5)  w24 → x[45,69]  y[ 2.5, 8.5]  在扫描管内
%  N3 支架标签   (2,34) north west w18 → x[ 2,20] y[25,34]
%  N4 针尖标签   (58,28) east     w13 → x[45,58] y[25.5,30.5]
%  N5 间距标签   (70,23) west     w30 → x[70,100] y[18.5,27.5]
%  N6 探测器标签 (80,73) south    w32 → x[64,96] y[73,81]
%  N7 激光标签   (2,68) south west w16 → x[ 2,18] y[68,77]
%  N8 反馈标签   中心(113,39) w32 → x[97,129] y[36,42]  在反馈盒内
% 两两求交核算（只列有风险的组合，均已排除）：
%  N3∩N4：x [2,20]∩[45,58]=∅；N4∩N5：x [45,58]∩[70,100]=∅；
%  N5∩反馈盒：y [18.5,27.5]∩[32,46]=∅；N5∩样品矩形：y 18.5>17；
%  N6∩探测器矩形：y 73>72；N7∩入射光束：光束在 x≤20 处 y≥57.7<68；
%  N1/N2 均在矩形内部；N8 在反馈盒内部（x 97..129 ⊂ 96..130，y 36..42 ⊂ 32..46）。
% 最大 x 坐标 = 130（反馈盒右边），最右节点文字到 x=129，故图宽 128 mm < 132 mm（\onecolwidth）。
% 控制线走 x=124（y 32→5.5），探测器信号线走 x=114（y 65→46），二者 x 不交，且都不与 N5 相交。
% ---- 样品与扫描管 ----
\fill[black!12] (30,10) rectangle (84,17);
\draw[axisline] (30,10) rectangle (84,17);
\draw[mainline] (30,17)--(84,17);
\node[align=center, text width=48mm] at (57,13.5) {样品（导体或绝缘体）};
\draw[axisline] (44,2) rectangle (70,9);
\node[align=center, text width=24mm] at (57,5.5) {xy 扫描管\\（压电陶瓷）};
% ---- 悬臂、支架与针尖 ----
\draw[axisline] (12,36) rectangle (26,44);
\node[align=left, text width=18mm, anchor=north west] at (2,34) {微悬臂\\与支架};
\draw[mainline] (24,40)--(64,40);
\fill[black!35] (61,40)--(67,40)--(64,19)--cycle;
\node[align=center, text width=13mm, anchor=east] at (58,28) {针尖};
\draw[thinline,dashed] (59,28)--(62.5,29.5);
% ---- 激光光路（入射 6,66 → 48,40.5；反射 48,40.5 → 77.4,58）----
\draw[mainline] (6,66)--(48,40.5);
\draw[mainline] (48,40.5)--(58,46.4);
\draw[mainline,->] (58,46.4)--(64,50.0);
\draw[mainline] (64,50.0)--(77.4,58);
\node[align=left, text width=16mm, anchor=south west] at (2,68) {激光\\二极管};
% ---- 四象限探测器 ----
\draw[axisline] (68,58) rectangle (92,72);
\draw[thinline] (80,58)--(80,72);
\draw[thinline] (68,65)--(92,65);
\node[align=center, text width=32mm, anchor=south] at (80,73) {四象限\\光电探测器};
% ---- 反馈控制 ----
\draw[axisline] (96,32) rectangle (130,46);
\node[align=center, text width=32mm] at (113,39) {反馈控制器与\\$z$ 压电陶瓷};
\draw[axisline] (92,65)--(114,65)--(114,46);
\draw[axisline] (124,32)--(124,5.5)--(70,5.5);
% ---- 针尖--样品间距标注 ----
\node[align=left, text width=30mm, anchor=west] at (70,23) {针尖--样品间距 $d$\\（示意，已放大）};
\draw[thinline,dashed] (64,19)--(70,23);
\end{tikzpicture}
\caption{光束偏转法 AFM 的原理示意。读者应看三条通路：悬臂偏转改变反射光方向，光斑在
四象限探测器上的位移给出偏转量；探测器信号送入反馈控制器，控制器驱动 $z$ 压电陶瓷使
针尖--样品间距回到设定值；扫描管在 $xy$ 平面内移动样品，记录的 $z$ 位移即形貌通道。
图中针尖--样品间距为便于作图已放大，实际成像时间距在纳米量级；悬臂弹性常数 $k$ 把
偏转量换算为力（$F=k\Delta z$）。本图为按工作原理重画的示意图，不代表具体仪器尺寸。
\schemamark\figsrc{据 S2 图 12-21（b），书内 p.913--914 / PDF p.399--400 重绘。}}
\label{fig:afm-setup}
\end{figure}

\subsection{成像模式与力--距离曲线}
\label{sec:afm-modes}

AFM 主要有三种基本工作模式：接触模式（contact mode）、轻敲模式（tapping mode）与非接触
模式（non-contact mode）。三者的共同点是：针尖在样品表面扫描时，样品--表面相互作用力的
改变使悬臂的偏转或振动模式改变，通过检测这一改变来反映表面的形貌或力学特征；差别在于
\textbf{把哪一个量送进反馈回路}，由此决定了工作距离、横向力大小与适用的样品范围
\src{S2 §12.3.2，书内 p.914 / PDF p.400}。

\subsubsection{接触模式：以静态偏转为反馈量}

接触模式又称静态 AFM，直接以悬臂的偏转作为力的度量。针尖与样品距离很远时悬臂处于零偏转
位置；针尖靠近或远离表面时针尖--表面作用力改变，悬臂偏转程度随之改变。偏转程度作为反馈
参数送给反馈控制器，控制器通过反馈运算控制 $z$ 方向压电陶瓷向上或向下运动，以保持恒定的
针尖--样品相互作用力；记录 $z$ 压电陶瓷的运动信号即反映样品形貌
\src{S2 §12.3.2，书内 p.914--915 / PDF p.400--401}。

静态 AFM 扫描时针尖通常与表面直接接触：这样能获得较好的分辨率，但容易损坏样品。采用较软
的悬臂（使其末端原子键能小于表面原子键能）可以减少对表面的损伤，但针尖靠近表面时又更
容易发生不受控制的跳变接触（jump-to-contact），导致难以稳定成像。正是为了回避这一对矛盾，
AFM 才发展出动态模式\src{S2 §12.3.2，书内 p.915 / PDF p.401}。这里要注意一个判断顺序：
"接触模式分辨率最高"这句话只在硬而平整的样品上成立；对软样品，横向剪切力造成的形变会同时
降低分辨率并改变高度读数。

\subsubsection{轻敲模式：以振幅（或相位）为反馈量}

轻敲模式是目前大气环境或液体中 AFM 最常用的工作模式。它把微悬臂以固定频率和振幅驱动振动，
振幅一般相对较大（\qtyrange{20}{100}{nm}），反馈信号取悬臂的振动\textbf{振幅或相位}。
轻敲模式对样品的损坏小，适用于不同的材料，具微米级扫描范围，对表面起伏程度大的样品适应性
好，且能在大气和液体中工作；代价是响应时间比较长、扫描速度较慢
\src{S2 §12.3.2，书内 p.915 / PDF p.401}。

轻敲模式之所以能同时给出"形貌"与"力学"两套信息，是因为它有两个可记录通道：振幅通道用于
反馈（维持设定振幅），相位通道是自由量。相位滞后反映针尖与样品相互作用的能量耗散与局部
黏弹性差异，因此\textbf{相位像常常比高度像更能区分同高度的不同组分}（例如载体上的金属
颗粒与载体本身）。但相位像不能按高度解读，这是本节后面单独列出的误判之一。

\subsubsection{非接触模式：以共振频率偏移为反馈量}

非接触式 AFM（non-contact AFM，NC-AFM）常用于真空环境。按检测信号与反馈参数，它又分为
振幅调制模式（amplitude-modulation mode，AM mode）与频率调制模式（frequency-modulation
mode，FM mode）。后者（FM-AFM）是动态 AFM 中较易取得原子级空间分辨率的工作模式：微悬臂
振幅一般小于 \qty{5}{nm}，反馈参数是悬臂共振频率的偏移，而共振频率的变化直接反映了
针尖--样品表面间的相互作用力梯度。近年来 Giessibl 发明的 qPlus 型传感器经不断改进，空间
分辨率达到亚原子级，已成功实现单分子内的原子分辨，是目前 AFM 技术与应用研究的前沿
\src{S2 §12.3.2，书内 p.915 / PDF p.401}。

动态 AFM 能够有效消除不稳定的针尖--样品接触，还能减小偏转测量中的噪声影响，因此目前
实际操作的多数是动态 AFM 模式（轻敲与非接触）
\src{S2 §12.3.2，书内 p.915 / PDF p.401}。三种模式的对比见表~\ref{tab:afm-modes}。

\begin{table}[htbp]
\centering
\caption{AFM 三种基本工作模式的对比}
\label{tab:afm-modes}
{\srcsize
\begin{tabularx}{\columnwidth}{@{}lXXX@{}}
\toprule
对比项 & 接触模式（静态 AFM） & 轻敲模式 & 非接触模式 \\
\midrule
反馈量 & 悬臂静态偏转（等价于 $z$ 位移） & 振动振幅（相位为附加通道） & 共振频率偏移 $\Delta f$（FM）或振幅（AM） \\
横向力 & 大：针尖与表面持续接触并剪切 & 小：间歇接触 & 几乎为零：不接触 \\
悬臂振幅 & 无（静态工作） & \qtyrange{20}{100}{nm} & 一般 $<\qty{5}{nm}$；qPlus 可达 \qty{100}{pm} \\
常用环境 & 大气、液体（须抑制毛细力） & 大气、液体（最常用） & 超高真空（FM 为主） \\
适用样品 & 硬、平整、与基底结合牢固 & 软样品、起伏大、松散颗粒 & 平整单晶，绝缘体亦可 \\
横向分辨率 & 理想条件下可达原子级 & 通常 \qtyrange{1}{10}{nm} & 原子级；FM-qPlus 达亚原子级 \\
主要优点 & 结构简单、速度快、可同步测摩擦 & 损伤小、稳定性好、环境适应性好 & 不损伤表面，可获得短程力主导衬度 \\
常见假象 & 划痕与拖拽，软样品被压陷 & 相位像被误当形貌，探针污染后衬度失真 & 长程力背景未扣除，振幅设定不当 \\
\bottomrule
\end{tabularx}}
\src{S2 §12.3.2，书内 p.914--915 / PDF p.400--401}
\end{table}

\subsubsection{动态 AFM 的谐振子模型：频率偏移为什么反映力梯度}

理解动态模式的反馈量，需要一个最小的力学模型。在动态 AFM 中微悬臂以其共振频率振动，
整个微悬臂系统可近似为一个处于外界激励力下的弹簧谐振子；针尖--样品作用力
$F_{\mathrm{ts}}$ 可看作对悬臂简谐振动的微扰，此时微悬臂是受迫阻尼振动。由此可得到微悬臂
的三个重要性质参数：弹性常数 $k$、共振频率 $\omega$ 与品质因数 $Q$
\src{S2 §12.4，书内 p.915 / PDF p.401}。

\begin{principle}
\textbf{第一步，针尖远离表面}（$F_{\mathrm{ts}}=0$）：微悬臂只受外界激励力 $F_{\mathrm{drive}}$
与阻尼作用 $F_{\mathrm{frict}}$。取余弦激励 $F_{\mathrm{drive}}=F_0\cos\omega t$（$F_0$ 为
激励振幅，$\omega$ 为激励频率），阻尼摩擦力与速度成正比 $F_{\mathrm{frict}}=m\gamma\,\dd z/\dd t$
（$m$ 为振子质量，$\gamma$ 为阻尼系数），振子离平衡位置的距离为 $z$，则受迫振动方程为
\src{S2 §12.4.1，书内 p.916 / PDF p.402}：
\end{principle}

\begin{equation}\label{eq:afm-driven}
m\frac{\dd^{2}z}{\dd t^{2}}=-kz-m\gamma\frac{\dd z}{\dd t}+F_0\cos\omega t
\end{equation}

式中 $k$ 为弹性常数。对自由谐振子，弹性常数与自然共振频率 $\omega_0$ 之间存在关系
$\omega_0=\sqrt{k/m}$；品质因数定义为 $Q=\omega_0/\gamma$：

\begin{equation}\label{eq:afm-q}
\omega_0=\sqrt{\frac{k}{m}},\qquad Q=\frac{\omega_0}{\gamma}
\end{equation}

给定参数后式~\eqref{eq:afm-driven} 有解析解，其解包含初始激励导致的瞬态项与稳定项；
经过足够长时间衰减之后，振动可近似为 $z=A\cos(\omega t-\varphi)$，振幅 $A$ 与相位偏移
$\varphi$ 都是激励频率的函数\src{S2 §12.4.1，书内 p.916 / PDF p.402}：

\begin{equation}\label{eq:afm-amp}
A(\omega)=\frac{F_0/m}{\sqrt{\left(\omega^{2}-\omega_0^{2}\right)^{2}+\left(\omega\omega_0/Q\right)^{2}}}
\end{equation}

\begin{equation}\label{eq:afm-phase}
\tan\varphi=\frac{\omega\omega_0/Q}{\omega_0^{2}-\omega^{2}}
\end{equation}

即外界激励下悬臂的运动仍是余弦函数，振动频率等于激励频率 $\omega$，但相对激励力有一个相位
延迟 $\varphi$。式~\eqref{eq:afm-amp} 与式~\eqref{eq:afm-phase} 是后面所有动态模式读数的
基础：\textbf{振幅在共振时取极大，相位在共振时过 $90$°，因此"把激励频率定在何处"直接
决定了反馈量对力变化的灵敏度}。

\textbf{第二步，针尖靠近表面}，且先假设针尖--表面相互作用力在平衡位置附近随距离线性变化。
此线性近似下，力梯度可等效为一个新的弹簧弹性常数
$k_{\mathrm{ts}}=-\left.\partial F_{\mathrm{ts}}/\partial z\right|_{z=0}$，运动方程中的
$k$ 变为 $k+k_{\mathrm{ts}}$，由此定义有效弹性力常数 $k_{\mathrm{eff}}=k+k_{\mathrm{ts}}$
与有效共振频率 $\omega_{\mathrm{eff}}=\sqrt{k_{\mathrm{eff}}/m}$。针尖--表面作用力一般远
小于悬臂回复力，即 $\left|k_{\mathrm{ts}}\right|\ll k$，于是共振频率相对自然共振频率的
偏移可近似为\src{S2 §12.4.2，书内 p.917 / PDF p.403}：

\begin{equation}\label{eq:afm-domega-lin}
\Delta\omega=\omega_{\mathrm{eff}}-\omega_0\approx\omega_0\frac{k_{\mathrm{ts}}}{2k}
=-\frac{\omega_0}{2k}\left.\frac{\partial F_{\mathrm{ts}}}{\partial z}\right|_{z=0}
\end{equation}

式~\eqref{eq:afm-domega-lin} 的结论很干净：\textbf{共振频率偏移正比于针尖--表面相互作用力
的梯度}。对正的作用力梯度（$\partial F_{\mathrm{ts}}/\partial z>0$，即吸引随靠近而增强），
$\Delta\omega<0$，共振频率向低频移动；对负的作用力梯度，共振频率向高频移动。注意源书同时
提醒：实际条件下若振动幅度较大，线性假设失效，因此该方法不能直接用于定量分析，但谐振子
模型给出的解析结果对理解动态 AFM 工作原理仍然非常有价值
\src{S2 §12.4.2，书内 p.917 / PDF p.403}。

\textbf{第三步，回到实际工作状态}。振幅调制 AFM 的振幅很大（典型约 \qty{50}{nm}），线性
近似已经失效：激励频率固定且往往接近或等于自然共振频率，针尖在一个振动周期内可深入排斥力
区与吸引力区。但实验与数值计算都证明，\textbf{悬臂振动幅度随针尖--样品间距的减小而线性
下降}，因此以振幅作为线性反馈信号来控制针尖与样品的距离依然十分有效，这也是振幅调制 AFM
被广泛使用的一个原因\src{S2 §12.4.3，书内 p.917 / PDF p.403}。

频率调制 AFM 则始终把激励频率保持在实际共振频率附近（$\omega\approx\omega_0$）。在针尖--
表面间的非线性作用下，悬臂频率偏移为\src{S2 §12.4.3，书内 p.917--918 / PDF p.403--404}：

\begin{equation}\label{eq:afm-domega-fm}
\Delta\omega=-\frac{\omega_0}{A^{2}k}\left\langle F_{\mathrm{ts}}\,z\right\rangle
\end{equation}

其中 $\left\langle F_{\mathrm{ts}}\,z\right\rangle$ 定义为一个振动周期内的平均：

\begin{equation}\label{eq:afm-fz}
\left\langle F_{\mathrm{ts}}\,z\right\rangle
\equiv\frac{1}{T}\int_{0}^{T}F_{\mathrm{ts}}\left(d+A\sin\omega t\right)A\sin\omega t\,\dd t,
\qquad T=\frac{1}{\omega_0}
\end{equation}

可以看出，频率偏移正比于 $\left\langle F_{\mathrm{ts}}\,z\right\rangle$，即一个振动周期 $T$
内针尖--样品作用力与针尖--样品距离乘积的平均。特别地，当悬臂振幅很小时（约 \qty{100}{pm}），
可用线性近似结果，$\Delta\omega$ 简化为式~\eqref{eq:afm-domega-lin}，即频率偏移正比于
针尖--表面相互作用力梯度\src{S2 §12.4.3，书内 p.918 / PDF p.404}。

\tocheck{未对照原页面图像：式~\eqref{eq:afm-domega-fm} 的分母，源文 OCR 作 $A^{2}k$；
按 Giessibl 原始推导，相同定义的 $\left\langle F_{\mathrm{ts}}z\right\rangle$ 前应为
$2kA^{2}$。此处暂按源文形式排，待核原页面后再定是否补因子 2。}

这一小振幅条件正是 qPlus 的工作点：其石英音叉弹性常数很高（$k=\qty{1800}{N/m}$），因此
可以稳定地以 \qty{100}{pm} 量级的振幅振动，刚好工作在线性近似区域；又因为短程力梯度很大，
可有效消除长程力背景对成像的影响，获得短程力主导的成像衬度，使 AFM 超越了 STM 的分辨，
达到亚原子分辨层次\src{S2 §12.4.3，书内 p.918 / PDF p.404}。

\subsubsection{力--距离曲线能给出什么}
\label{sec:afm-forcecurve}

前一小节说明的是悬臂本身；把针尖停在某一点、让 $z$ 压电陶瓷走一条"逼近--回撤"的曲线，
同时记录悬臂偏转，就得到\textbf{力--距离曲线}（force--distance curve）；在每个像素都做一次
即成\textbf{力学谱}（force spectroscopy / 空间分辨力学谱）。

\begin{measure}
测力曲线时必须同时记录并给出的条件有：悬臂弹性常数 $k$ 及其标定方法、针尖半径 $R$、
$z$ 压电陶瓷的位移量程与速度、环境（真空/干燥气/大气/液体）与温度。曲线横轴有两种取法：
一是 $z$ 压电位移，二是针尖--样品间距 $d$；两者在接触区相差悬臂偏转量 $\Delta z$，
换算关系就是胡克定律\src{S2 §12.3.2，书内 p.914 / PDF p.400}：

\begin{equation}\label{eq:afm-hooke}
F=k\,\Delta z
\end{equation}

因此未给出 $k$ 的力曲线只能做相对比较，不能给出力的绝对值。
\end{measure}

一条完整的力--距离曲线（图~\ref{fig:afm-forcecurve}）可直接读出三个量：

\begin{itemize}
\item \textbf{黏附力（最大吸引力）}：回撤支上针尖脱离表面处的跳离点，对应偏转的极大值
$\Delta z_{\max}$，由式~\eqref{eq:afm-hooke} 得 $F_{\mathrm{ad}}=k\,\Delta z_{\max}$。
它反映针尖与表面之间所有吸引作用之和。
\item \textbf{刚度（弹性模量）}：接触区曲线斜率即有效刚度 $\dd F/\dd d$；在弹性接触模型下
可由斜率反推材料的有效弹性模量，但需要已知针尖半径 $R$ 与接触几何。
\item \textbf{过零距离与短程力特征}：总力为零处对应近似成键距离；最大吸引的位置与大小、
以及短程力的梯度，是与元素种类直接相关的量（§\ref{sec:afm-app} 中单原子区分的依据）。
\end{itemize}

\begin{figure}[htbp]
\centering
\begin{tikzpicture}
\begin{axis}[
  width=\onecolwidth, height=6.2cm,
  xlabel={针尖--样品间距 $d$ / nm},
  ylabel={相互作用力（模型单位）},
  xmin=0.15, xmax=1.5, ymin=-0.22, ymax=0.5,
  xtick={0.2,0.4,0.6,0.8,1.0,1.2,1.4},
  ytick={-0.1,0,0.1,0.2,0.3,0.4,0.5},
  % 图例必须放在轴内：实测 legend pos=outer north east 会把整图撑到 483 pt（170 mm），
  % 超出 165 mm 版心 13.9 pt，成为 overfull 的唯一根因。轴内放置后整图约 133 mm。
  legend pos=north east, legend cell align=left,
]
\addplot[domain=0.185:1.5,samples=400,mainline,black]
  {0.34*exp(-(x-0.26)/0.075)-0.34/(1+((x-0.26)/0.12)^2)};
\addlegendentry{总力 $F_{\mathrm{ts}}$}
\addplot[domain=0.185:1.5,samples=400,thinline,black,dashed]
  {-0.34/(1+((x-0.26)/0.12)^2)};
\addlegendentry{长程吸引（范德华）}
\addplot[domain=0.185:1.5,samples=400,thinline,black,dotted]
  {0.34*exp(-(x-0.26)/0.075)};
\addlegendentry{短程排斥（Pauli）}
\draw[thinline,dashed] (axis cs:0.26,-0.22)--(axis cs:0.26,0.5);
% 注释节点坐标区间静态核算（数据域 x 0.15--1.5、y -0.22--0.5，轴盒约 117 mm x 50 mm）：
%  N-a 零力点  (0.30,0.30) south west w40mm → x[0.30,0.76] y[0.30,0.47]
%  N-b 图例（轴内 north east）              → x[0.98,1.50] y[0.33,0.50]
%  N-c 吸引区  (0.85,0.16) west w22mm      → x[0.85,1.10] y[0.13,0.19]
%  N-d 黏附力  (0.41,-0.17) west w46mm     → x[0.41,0.94] y[-0.20,-0.14]
% 两两求交：N-a∩N-b x 不交(0.76<0.98)；N-c∩N-b y 不交(0.19<0.33)；N-d 距曲线最近约 1.2 mm；
% N-d 下边界 -0.20 > ymin -0.22，不会被轴裁掉。
\node[align=left, text width=40mm, anchor=south west] at (axis cs:0.30,0.30)
  {零力点 $\approx$ 成键距离（$d\approx\qty{0.26}{nm}$）};
\draw[thinline,dashed] (axis cs:0.335,-0.118)--(axis cs:0.335,-0.17);
\draw[thinline,dashed] (axis cs:0.335,-0.17)--(axis cs:0.40,-0.17);
\node[align=left, text width=46mm, anchor=west] at (axis cs:0.41,-0.17)
  {最大吸引即黏附力的来源};
\node[align=left, text width=22mm, anchor=west] at (axis cs:0.85,0.16) {吸引区};
\end{axis}
\end{tikzpicture}
\caption{针尖--样品相互作用力随间距的变化（示意模型曲线）。读者应看三处：$d$ 大于零力点时
总力为吸引且随间距增大缓慢趋零（长程范德华项）；$d\approx\qty{0.26}{nm}$ 处总力过零
（近似成键距离）；$d$ 小于该值时总力转为排斥并急剧上升（Pauli 排斥，曲线在纵轴顶端被截断）。
虚线为长程吸引分量，点线为短程排斥分量，实线为两者之和。横轴为针尖--样品间距；实测力曲线
的横轴通常是 $z$ 压电位移，两者在接触区相差悬臂偏转量。本图为按长程吸引与短程排斥两项叠加
计算的示意曲线，纵轴为模型单位而非实测力值，不代表任何具体针尖--样品组合。
\modelmark\figsrc{据 S2 图 12-21（a），书内 p.913--914 / PDF p.399--400 改绘并简化。}}
\label{fig:afm-forcecurve}
\end{figure}

为什么力曲线比形貌像更接近"材料性质"？关键在两种数据的性质不同。形貌像的高度是\textbf{反馈
回路维持设定点所需的 $z$ 位移}：它把针尖--样品作用力的空间变化压成一条等值面，等值面上
任意两点只要"力相同"就显示为同一高度，几何起伏、电子态差异、吸附层与形变被混在一起，
无法分离。力曲线是逐点的 $F(d)$，直接给出作用力的\textbf{绝对值与距离依赖}：最大吸引、
过零距离、短程力梯度都由针尖原子与表面原子的种类和成键决定。因此，力学谱可以在不同针尖、
不同实验之间通过归一化比较（§\ref{sec:afm-app} 的 Sn/Si、Pb/Si 比值就是例子），而形貌像
的亮度只有相对意义。定量化弹性模量还需要额外条件（已知 $R$、标定过的 $k$、接触模型适用），
这一点不能省略。

\subsubsection{针尖卷积：图像是针尖与样品的形态卷积}
\label{sec:afm-conv}

AFM 的横向分辨率不由衍射极限决定，而由\textbf{针尖末端的形状与半径 $R$} 决定：测得的图像
是针尖形态与样品形态的"叠加"（几何上即两者的 Minkowski 和），凡比针尖更尖、更陡的结构都
会被针尖"描圆"。

\begin{principle}
设针尖末端为半径 $R$ 的球，样品上有一个半径为 $r$ 的孤立球形颗粒（$r\ll R$），扫描时
反馈维持恒定的针尖高度。颗粒中心正上方测得的高度为 $2r$（即颗粒直径，此时正确）；在偏离
中心 $x$ 处，针尖顶点高度为 $h(x)=r-R+\sqrt{(R+r)^{2}-x^{2}}$。以半高（$h=r$）度量宽度，
得表观半高宽
\end{principle}

\begin{equation}\label{eq:tipconv}
w_{\mathrm{app}}=2\sqrt{2Rr+r^{2}}\;\approx\;2\sqrt{2Rr}\qquad(r\ll R),
\qquad
\frac{w_{\mathrm{app}}}{2r}\approx\sqrt{\frac{2R}{r}}
\end{equation}

即\textbf{表观横向尺寸随 $\sqrt{R}$ 增长、随真实尺寸的平方根减小}。取 $R=\qty{10}{nm}$、
$r=\qty{1}{nm}$，表观半高宽约 \qty{9.2}{nm}，而真实宽度只有 \qty{2}{nm}——横向尺寸被高估
约 4.6 倍。

\supp{式~\eqref{eq:tipconv} 是按球形针尖与球形特征做几何卷积的标准结果，源书未给解析式；
它假设 $r\ll R$、以半高宽度量、忽略形变与吸附层。若按颗粒底部（基底高度）度量宽度，结果为
$4\sqrt{Rr}$，数值更大。}

\begin{pitfall}
\textbf{由 AFM 像直接报粒径分布是常见错误。}识别针尖卷积的四个可用判据：同一特征在不同
扫描方向或不同针尖下宽度不同；图像中出现重复出现的相同形状（多针尖/双针尖假象）；台阶边缘
被圆弧化，圆角半径与针尖半径同量级；高宽比大的结构被"抹平"。正确做法是用已知尺寸的标样
（金胶体、光栅）标定针尖，或换用更尖的针尖（碳纳米管针尖、CO 修饰针尖、qPlus 传感器）。
在针尖未标定的情况下，AFM 给出的横向尺寸只能当作\textbf{上界}。
\end{pitfall}

\subsubsection{误判与边界}
\label{sec:afm-pitfall}

\begin{pitfall}
（1）\textbf{把 AFM 高度当作真实几何高度。}形貌通道给出的是维持设定点所需的 $z$ 位移，
其中包含吸附层（水膜、羟基、残留碳）与针尖/样品形变：大气环境下氧化物表面常覆盖一层水膜
（\pedset{厚度取 \qtyrange{0.2}{0.5}{nm} 作为教学讨论量}），台阶高度会整体偏大；软样品在
接触模式下被压陷，高度又偏小。判据是改变设定点或环境后高度是否随之改变——若改变，则读出值
不是几何高度。\\
（2）\textbf{忽略针尖卷积，用图像尺寸当真实尺寸。}见式~\eqref{eq:tipconv}：常见针尖
（$R$ 为几纳米至几十纳米）会把纳米颗粒的横向尺寸放大数倍。颗粒计数得到的粒径分布、岛宽度、
孔口尺寸都受此影响。\\
（3）\textbf{把轻敲相位像当成纯形貌像。}相位滞后反映能量耗散与局部黏弹性，与高度不是同一
物理量；相位像与高度像经常不相关甚至反相。用相位像的高低谈"几何高低"是错的，但用相位像
区分同高度的不同组分（金属颗粒与载体）是恰当的用法。\\
（4）\textbf{忽略空气中吸附水膜对力曲线的影响。}针尖接近带水膜的表面时，针尖与水膜之间
形成毛细弯月面，附加一个吸引力：结果是远距离就出现吸引、跳变接触提前、黏附力偏大，
于是由曲线反推的弹性模量与最大吸引随之失真。判断方法是在干燥气氛、真空或液体中复测同一
位置，比较跳离点位置与黏附力大小。
\end{pitfall}

\subsection{原子力显微镜在催化研究中的应用}
\label{sec:afm-app}

AFM 在催化研究中的价值集中在一点：\textbf{它能在绝缘的氧化物、分子筛载体上做原子级结构
表征，并同时给出力（力学）信息}。以下五类应用按"看结构—认元素—看成键—看反应—看界面"
的顺序展开。

\subsubsection{氧化物表面结构表征}

氧化物是催化剂的常见组成部分：可作为载体负载金属纳米粒子，形成金属--氧化物界面的双活性
中心，也可作为吸光载体用于光催化。因此原子尺度的氧化物表面结构表征，对理解氧化物参与的
催化反应过程有重要意义。由于导电性较差，氧化物表面经常难以直接用 STM 观测；AFM 在绝缘体
表面也能正常工作，而且还能对针尖--表面间作用力做详细的力谱表征，为表面结构及表面吸附质
物种的判别提供更多信息。目前 AFM 已在 \ce{CeO2}(111) 和 \ce{TiO2} 等氧化物表面实现较为
清楚的原子分辨，其他重要氧化物催化剂表面如 $\alpha$-\ce{Al2O3} 也可用 AFM 实现原子级
表征\src{S2 §12.5.1，书内 p.918 / PDF p.404}。

一个特别有用的做法是\textbf{AFM 与 STM 联用}：采用导电金属针尖，在记录 AFM 信号的同时
记录隧穿电流的变化，可同时得到两套衬度。锐钛矿型 \ce{TiO2}(101) 表面就是典型例子——通过
详细对比 STM 与 AFM 图像的衬度特征，并结合第一性原理计算结果，Perez 与 Custance 等成功
识别了表面吸附水分子、次表层氧缺陷以及表面羟基。利用空间分辨的力学谱还可判定：AFM 像中的
亮度主要来自表层氧，而 STM 像中的亮点是次表层 Ti$_{5c}$ 位点。原因是 AFM 更多反映最表层
O 原子的信息，而同时记录的 STM 电流信号对表面局域电子态更敏感，因此能从电流图中看出次表层
Ti 位点；AFM 力学谱还能给出金红石型 \ce{TiO2}(110) 表面吸附金属原子与羟基的识别信息
\src{S2 §12.5.1，书内 p.918--919 / PDF p.404--405}。

\ce{CeO2}(111) 表面的氧缺陷研究进一步显示了"形貌 + 耗散"双通道的价值。\ce{CeO2}(111) 是
氧终端面，表层氧缺陷与次表层氧缺陷在 AFM 图像中表现出不同特征：表层氧缺陷表现为一个缺失
原子被周围六个突起的氧原子围绕；次表层氧缺陷则在表面没有原子缺失，但氧空位上方间距为两倍
晶格常数的三个氧原子在 AFM 图中表现突起，其内的三个相邻氧原子则向内稍微弛豫，这些实验观察
与第一性原理计算结果吻合。AFM 在记录拓扑图像时还可同时记录维持悬臂恒振幅振动的
\textbf{能量耗散信号}（dissipation signal），这种能量耗散由 AFM 针尖尖端原子与表面的相互
作用导致。Torbrugge 等通过对比 AFM 拓扑图像与能量损失信号图像发现，次表层氧空位上方三个
紧邻氧原子与针尖作用较强，其能量损失信号强于周围突起的氧原子，由此推测这些紧邻氧原子应该
比周围的氧原子具有更高的反应活性\src{S2 §12.5.1，书内 p.920 / PDF p.406}。

\tocheck{未对照原页面图像：源文在描述次表层氧缺陷形貌时标注为"图 12-22（i）"，与表层氧缺陷
同一编号，疑为（j）之误；此处只按文字描述转述，未按图号断言。}

\subsubsection{单原子尺度的元素区分}

单原子或单分子层次的化学分辨对 STM 和 AFM 一直是特别有挑战的难题，动态 AFM 的出现使这个
难题取得突破。2007 年 Custance 等利用频率调制 AFM（FM-AFM），在室温真空环境下成功实现了
Si--Sn--Pb 合金表面单个元素的化学区分
\src{S2 §12.5.2，书内 p.920 / PDF p.406}。

方法本身包含一个很值得学的归一化思想。针尖与表面之间的相互作用力有长程项（如范德华力）
与短程的化学成键作用力；他们发现虽然不同针尖会影响实验测得的针尖--表面原子力学曲线，
但\textbf{即便对于不同针尖，不同原子与针尖之间最大吸引力的比值是一个定值}。因此他们用同一
针尖测 Si、Pb、Sn 三种元素的力学曲线，并分别用该针尖在 Si 上的最大吸引力归一化，再由归一化
曲线求最大吸引力之比——这个比值与元素直接相关，例如 Sn/Si 为 $0.77$，Pb/Si 为 $0.59$。
通过这一最大力比值，就可以排除针尖因素的干扰，成功实现不同元素单个原子的区分
\src{S2 §12.5.2，书内 p.920 / PDF p.406}。

利用 AFM 力学测量技术，Gross 等还实现了识别 NaCl 表面不同荷电状态的金属原子：在同样的
针尖--样品间距下，与中性 Au$^{0}$ 原子相比，带负电的 Au$^{-}$ 与针尖的吸引力作用更大
\src{S2 §12.5.2，书内 p.920 / PDF p.406}。这两项工作合起来说明：AFM 不仅能获得氧化物负载
单原子催化剂的结构信息，还有望从原子层面理解金属--氧化物之间的电子转移规律。

\subsubsection{金属团簇内部与衬底的键合结构确认}

对负载金属纳米粒子，原子尺度的结构信息——金属内部成键及其与衬底的成键——对从根本上理解
纳米催化剂的结构--性能关系有重要作用，但长期未能取得直观的成像证据。2015 年 Giessibl 等
使用 qPlus-AFM 与 CO 分子修饰的针尖，对 Fe$_1$ 至 Fe$_{16}$ 团簇成像，成功实现了金属团簇
内部成键的观测。CO 修饰针尖极为清楚地反映了 Fe 团簇因表面吸附导致的电荷密度再分布，把表面
吸附原子的成像分辨率推到亚原子级别。作者再结合衬底的分辨图像确定成键位点：Fe 原子双聚体与
三聚体在表面的落位并不是广为人知的表面 fcc 位点，而更接近桥位（bridge）；而对于 Fe 单体与
Fe 四聚体以上的团簇，所有 Fe 都落在更稳定的表面 fcc 位\src{S2 §12.5.3，书内 p.921 /
PDF p.407}。

这个例子的方法论意义大于结论本身：\textbf{团簇尺寸改变落位}——小团簇落在低配位桥位，
大团簇回到密排的 fcc 位。这类信息直接影响对活性位结构的判断（例如金属--载体界面的配位
环境），而它只能由原子级成像给出。

\subsubsection{化学键的实空间成像与反应中间体}

qPlus-AFM 极高的空间分辨率使人可以直接观察共价键，甚至区分不同键级（bond order）的碳碳键；
而识别分子间相互作用力对揭示分子反应中间体、氢键的物理本质有重要意义。2013 年，国家纳米
中心的裘晓辉等在国际上首次实现氢键的直接观测：在 Cu(111) 表面沉积 8-羟基喹啉
（8-hydroxyquinoline，8-hq）分子，再用 qPlus-AFM 观测，发现分子聚集体之间形成了可观测的键，
结合 DFT 计算模拟以及成键位置和取向，确认为分子间氢键。由于 qPlus-AFM 对分子内成键的衬度
来源于针尖与样品原子间的 Pauli 排斥力，因此该衬度与电子密度的空间分布直接相关；分子间氢键
的成功观测，为氢键的共价键成分与 H 到 N/O 之间的电荷转移提供了直接的实验证据。作者进一步
观察到表面 Cu 吸附原子与脱氢 8-hq 之间形成金属--N 配位键，这为直接观测表面反应、表面吸附
与成键提供了可能\src{S2 §12.5.4，书内 p.922--923 / PDF p.408--409}。

反应中间体的直接观测是理解反应动力学的重要方法，但中间体活性高、表面浓度低，许多常规手段
难以获得其信息。qPlus-AFM 因为可以直接看到化学键，用于表面有机分子反应的观察时可以直接
分辨反应分子与多种反应产物。Crommie 等用 qPlus-AFM 实现了表面分子反应中间体的直接捕获与
观察：研究对象是 1,2-二（2-乙炔基苯基）乙炔在 Ag(100) 表面的交叉偶联与分子内环化反应，
作者不仅看到了反应分子和产物，还捕捉到了中间体。为了获得更多机理认识，作者把表面加热到
不同温度和时间，再用 qPlus-AFM 确定表面分子浓度，得出整个反应进程中产物与各中间体的变化。
他们发现，用经典密度泛函理论结合微观动力学无法模拟出实验观察到的中间体浓度变化，必须考虑
反应分子与表面之间的传能以及表面对不同中间体活化熵的影响，理论模拟结果才能与实验吻合。
结合理论计算与实验观察，作者得出结论：中间体在表面可以自由旋转和平动、活化熵较高，因此
参与下一步反应的速率较慢，从而可被实验观察到；而其他理论预测的中间体因为活化熵较低，
很快在反应中被消耗掉\src{S2 §12.5.4，书内 p.923--924 / PDF p.409--410}。

这一案例给催化研究者的提示是：\textbf{表面反应的选择性不只由活化焓决定，活化熵（即中间体
在表面的自由度）可以直接改变可观测的产物分布}，而 AFM 提供的是唯一能"点名"每个中间体的
实验通道。

\subsubsection{固--液界面的空间三维成像}

固--液界面的原子级观测是理解电化学与生命现象的关键。长期以来固体表面水分子的原子排列模型
都不特别明确；特别是水溶液--固体界面，由于双电层的存在，界面结构在空间上是局域变化的，
常规二维扫描很难得到界面结构在空间上的分布信息。由于 AFM 可以工作在较大的针尖--表面距离
范围，因此在三维空间上的直接观测可以解析固--液结构的完整图像
\src{S2 §12.5.5，书内 p.925 / PDF p.411}。

Fukuma 等用 AFM 直接观测了白云母片上的水分子排列，得到近表面层水分子排列图像，并结合力学
曲线获得了离固体表面垂直距离 \qty{0.78}{nm} 以内的水--固体界面三维图像。观测结果是分层
的：离白云母表面较远处，AFM 基本看不到明显衬度；针尖更加靠近表面后，$xy$ 平面截图才给出
结构。在 \qty{0.2}{mol/L} KCl 溶液中，三维 AFM 图像显示表面吸附的 K$^{+}$ 单层及其顶端
两层水合层，水合层约 \qty{0.3}{nm} 厚，随表面起伏排列；在 \qty{4}{mol/L} KCl 溶液中，
界面层分为表面约 \qty{2}{nm} 厚的有序液相层与体相溶液层；在约 \qty{5}{mol/L} NaCl 溶液中，
则得到原子级有序的白云母--溶液界面三维 AFM 图像
\src{S2 §12.5.5，书内 p.925 / PDF p.411}。

对催化而言，这类结果的意义在于把"界面水"从连续介质变成了可分层的结构：水合层厚度、
吸附离子单层与有序液相层的存在，直接决定电催化界面上的传质与双电层模型假设是否成立。

\subsubsection{读图实例：从形貌通道与耗散通道判断氧空位类型}

下面用 \ce{CeO2}(111) 氧缺陷的判别过程演示 AFM 数据的读法。观察与推断分开陈述。
形貌特征取自源文对图 12-22（i）至（o）的文字描述
\src{S2 §12.5.1，书内 p.919--920 / PDF p.405--406}。

\begin{obs}
在一张 \ce{CeO2}(111) 的原子分辨 AFM 图上（同时记录形貌通道与能量耗散通道）：
① 缺陷 A：一个亮点的位置空缺，其周围六个亮点仍按六方排列，并相对周围略微凸起，形成
"缺一个原子的六角环"；
② 缺陷 B：表面上没有亮点缺失，但空位正上方、间距为两倍晶格常数的三个亮点连成一组凸起，
该组内部的三个相邻亮点相对周围略微内凹（弛豫）；
③ 在能量耗散通道中，缺陷 B 处那三个紧邻亮点的耗散信号强于周围的凸起亮点；
④ 两类缺陷在形貌通道中的高度起伏都在亚埃量级（\pedset{本实例不给出具体高度数值，
以免被当作从原图量取的结果}）。
\end{obs}

\begin{infer}
由 ① 推断缺陷 A 是\textbf{表层氧空位}：表面直接缺失一个氧，周围氧向空位弛豫而显得凸起。
由 ② 推断缺陷 B 不能用表层空位解释（表面没有亮点缺失），只能来自表面以下的缺陷：
\textbf{次表层氧空位}通过其上方氧原子的弛豫改变了针尖--表面作用力，从而在形貌通道产生
衬度。由 ③ 进一步推断：那三个氧原子与针尖的相互作用更强，源文据此推测它们比周围的氧原子
具有更高的反应活性。把两个通道合起来还能得到一个一般结论：在 \ce{CeO2}(111) 上"AFM 亮"
对应最表层氧，而在 \ce{TiO2} 上"STM 亮"对应次表层 Ti$_{5c}$ 位点——不同通道的"亮点"
不是同一种东西，跨方法比较亮度前必须先确认衬度来源。
\end{infer}

推断的边界必须写明：如果只有形貌通道而没有耗散通道与第一性原理计算的对照，②③ 只能算
假设；形貌衬度还会随针尖状态（CO 修饰、钝化程度）整体改变甚至反相，因此判断缺陷类型时
必须报告针尖状态与设定点。

\begin{dataflow}
从原始数据到结论的链条（以上例为例）：① 探测器电压 $\rightarrow$ 悬臂偏转或振幅信号
$\rightarrow$ 扣除远离表面处的参考值后得到频移/振幅通道；② 反馈回路把该通道维持在设定点，
$z$ 压电位移即形貌通道，维持恒定振幅所需的激励功率即能量耗散通道；③ 由形貌通道的周期性
（晶格常数、缺陷间距）定标空间坐标；④ 由缺陷处的局部形貌特征 + 耗散衬度 + DFT 计算对比
判定缺陷类型；⑤ 在缺陷处取一条 $F(d)$ 力学谱，由最大吸引与短程力特征判断亮点对应哪种
原子。每一步都依赖上一步的定标：若第 ③ 步的横向定标被针尖卷积放大（式~\eqref{eq:tipconv}），
第 ④ 步"两倍晶格常数"这一判据就会失效。
\end{dataflow}

\subsubsection{互补与回查}
\label{sec:afm-crosscheck}

\begin{crosscheck}
AFM 的强项（绝缘体可用、给力与力学性质）与弱项（不给电子结构、横向尺寸受针尖卷积影响、
高度不等于几何高度）决定了它的回查路线：\\
\textbf{与 STM 联用}：同一针尖同时记录隧穿电流与力信号，可把"形貌亮"与"电子态亮"分开——
这是 \ce{TiO2}(101) 上区分表层氧与次表层 Ti 位点的关键，也适用于任何导电性刚好够用的
样品。\\
\textbf{与衍射/散射方法回查}：由 AFM 量出的晶格常数、台阶高度、缺陷间距应与 XRD/LEED
给出的体相或表面晶格一致；若偏差系统性偏大，先怀疑横向定标（针尖卷积）而不是样品。\\
\textbf{与第一性原理计算回查}：缺陷类型的判定几乎都依赖"实验形貌 + 计算模拟"的对照
（\ce{CeO2} 与 \ce{TiO2} 两个例子都是如此）；反过来，计算给出的力曲线又能预测不同元素
的最大吸引力之比，与力学谱实测比值比较。\\
\textbf{与方法边界回查}：AFM 看不到的元素种类、氧化态与体相结构，仍要靠 XPS、EXAFS、
化学吸附滴定补足；把 AFM 的形貌结论直接当"活性位归属"是越界。
\end{crosscheck}

\srcnote{本节（12.3—12.5）内容据 S2（下册）第 12 章，书内 p.913—925 / PDF p.399—411。
其中式~\eqref{eq:lj}、式~\eqref{eq:tipconv} 及毛细力讨论为按学科标准形式的补充，
已在正文中单独标注；标注 \tocheck{} 之处待对照原页面图像核验。}
